\setcounter{section}{0}

\Needspace{5\baselineskip}
\hypertarget{outlook-on-key-milestones-in-embodied-ai}{%
\section{Outlook on Key Milestones in Embodied AI}\label{outlook-on-key-milestones-in-embodied-ai}}

Following the advent of large language models, AI has developed rapidly in the digital world. Embodied AI extends it into the physical world, closing the loop between perception, reasoning, and action in reality and bringing an important transformation in the form of intelligence. The following is an outlook on several key future milestones.

\Needspace{5\baselineskip}
\hypertarget{i.-three-key-milestones-in-embodied-ai}{%
\subsection{I. Three Key Milestones in Embodied AI}\label{i.-three-key-milestones-in-embodied-ai}}

\Needspace{5\baselineskip}
\hypertarget{the-chatgpt-moment}{%
\subsubsection{(1) The ChatGPT Moment}\label{the-chatgpt-moment}}

This milestone means that embodied AI can generalize, understand the physical world, and reach human levels in physical-intelligence capabilities such as movement and manipulation. It can correctly complete specific physical-world primitive tasks, such as serving dishes or folding clothes, from language instructions. We introduce a ``physical Turing test'': humans can no longer determine from outcomes alone whether work was completed by a human or a robot.

\Needspace{5\baselineskip}
\hypertarget{the-physical-api}{%
\subsubsection{(2) The Physical API}\label{the-physical-api}}

In LLM development, the transition from chatbot to coding agent is a crucial change in form, marking the ability to execute complex engineering tasks. In embodied AI, it means robots can be invoked through physical APIs: a human or model only needs to describe the desired product or task, and embodied AI decomposes it through physical reasoning and executes actions step by step until the product is made or the task completed. This enables closed-loop automation of industrial and agricultural production and natural-science research.

\Needspace{5\baselineskip}
\hypertarget{physical-self-improvement}{%
\subsubsection{(3) Physical Self-Improvement}\label{physical-self-improvement}}

The fundamental difference between human and animal physical intelligence is humans' ability to invent and use tools. In LLM development, the transition from coding for engineering implementation to autonomous research and self-improving AI for AI is also a key milestone. For embodied AI, this means moving beyond a production executor told the product's exact form to a researcher ``given only an optimization objective and autonomously exploring what to create to achieve it.'' It can autonomously invent and use tools, manufacture and improve embodied AI itself, define interfaces between intelligence and the physical world, and treat physical objects as targets for closed-loop optimization, achieving physical intelligence's self-improvement and autonomous iteration.

\Needspace{5\baselineskip}
\hypertarget{ii.-from-the-chatgpt-moment-to-physical-apis-basic-forms-and-industry-applications}{%
\subsection{II. From the ChatGPT Moment to Physical APIs: Basic Forms and Industry Applications}\label{ii.-from-the-chatgpt-moment-to-physical-apis-basic-forms-and-industry-applications}}

In the future, we hope embodied robots can perform real-world work like humans and be widely deployed across industries.

\Needspace{5\baselineskip}
\hypertarget{everyday-life}{%
\subsubsection{(1) Everyday Life}\label{everyday-life}}

Embodied AI will undertake housework, care, companionship, and home repairs. Rather than merely an appliance executing fixed motions, it will understand family members' real needs, perceive the states of objects, older adults, children, and pets in complex environments, and autonomously plan within safety boundaries---for example, tidying rooms, preparing meals, caring for older adults, and handling unexpected events.

\Needspace{5\baselineskip}
\hypertarget{services}{%
\subsubsection{(2) Services}\label{services}}

Embodied AI will enter food service, hotels, logistics, healthcare and nursing, retail, and public services. Future robots will perform repetitive delivery, cleaning, transport, and guidance tasks and respond flexibly to customers' language, expressions, and behavior, making services more consistent and less costly while freeing human workers from intensive, low-creativity labor.

\Needspace{5\baselineskip}
\hypertarget{industry}{%
\subsubsection{(3) Industry}\label{industry}}

Embodied AI will enable unmanned industrial production. Large models can translate future human requirements expressed in natural language or similar forms into executable production plans, while robots autonomously perform the entire physical process. They can operate assembly-line machinery, monitor production, and respond promptly to emergencies.

\Needspace{5\baselineskip}
\hypertarget{agriculture}{%
\subsubsection{(4) Agriculture}\label{agriculture}}

Embodied AI will move agriculture from reliance on experience and manual labor toward precision, automation, and data-driven operation. Robots can sow, weed, fertilize, harvest, identify pests and diseases, monitor soil, and sort crops, adjusting strategies in real time to weather, soil, moisture, and crop growth. This improves yields, reduces waste, and alleviates agricultural labor shortages.

\Needspace{5\baselineskip}
\hypertarget{iii.-embodied-ai-self-improvement-advancing-intelligence-itself}{%
\subsection{III. Embodied AI Self-Improvement: Advancing Intelligence Itself}\label{iii.-embodied-ai-self-improvement-advancing-intelligence-itself}}

Some cognitive theories hold that human intelligence did not arise solely through abstract symbolic reasoning but gradually developed through prolonged environmental interaction. Infants build an understanding of objects, space, causality, and others' behavior by seeing, hearing, and actively interacting with the world. If this theory is correct, intelligence with all human capabilities requires an embodied ``perception--reasoning--action'' loop: not only observing and understanding the world but actively acting on it and learning autonomously. Embodied AI can improve intelligence through data, models, compute, and other dimensions, enabling recursive self-improvement in the physical world.

\Needspace{5\baselineskip}
\hypertarget{mechanisms-of-embodied-ai-self-improvement-expanding-resources-and-enabling-autonomous-learning}{%
\subsubsection{(1) Mechanisms of Embodied AI Self-Improvement: Expanding Resources and Enabling Autonomous Learning}\label{mechanisms-of-embodied-ai-self-improvement-expanding-resources-and-enabling-autonomous-learning}}

\Needspace{4\baselineskip}
\begin{enumerate}
\def\labelenumi{\arabic{enumi}.}
\tightlist
\item
  Accelerating physical-world iteration to support continuous compute expansion
\end{enumerate}

Intelligence depends on expanding compute, but human production capacity currently limits this. Embodied AI can accelerate physical-world iteration through closed loops, continuously expanding available compute by increasing scale and optimizing processes, preventing compute from constraining intelligence.

\Needspace{4\baselineskip}
\begin{enumerate}
\def\labelenumi{\arabic{enumi}.}
\setcounter{enumi}{1}
\tightlist
\item
  Overcoming data limitations through active acquisition
\end{enumerate}

Human Internet data is limited in scale, modalities, and scope. Depending on humans to put data into computers and ``feed'' models may ultimately be unsustainable. Once embodied AI can interact autonomously with its environment, it can actively acquire needed real-world data (primarily visual, auditory, tactile, and sensor data in everyday environments) and verify its judgments. As its range of action grows, the breadth of accessible data distributions expands, theoretically providing indefinitely scalable data. Acting to obtain missing data is an important part of autonomous learning.

\Needspace{4\baselineskip}
\begin{enumerate}
\def\labelenumi{\arabic{enumi}.}
\setcounter{enumi}{2}
\tightlist
\item
  Expanding verification to form physical-world optimization loops
\end{enumerate}

LLM improvements in mathematics, coding, and agent tasks mainly rely on reinforcement learning with verifiable rewards. Embodied AI can extend verification into reality, making tasks requiring real-world interaction into closed loops that models can handle, enabling RL credit assignment and optimization.

\Needspace{4\baselineskip}
\begin{enumerate}
\def\labelenumi{\arabic{enumi}.}
\setcounter{enumi}{3}
\tightlist
\item
  Closing the loop of autonomous learning
\end{enumerate}

Autonomous learning means actively obtaining data and feedback through environmental interaction and using them to update understanding. In embodied AI, or next-generation models combining it with large models, the former includes not only state transitions and actions for deep learning but also verification signals for reinforcement learning; the latter can combine methods that optimize models from given data or rewards, such as AI for AI. Together they close an autonomous-learning loop, with compute expansion providing fundamental support, echoing the search and learning discussed in ``The Bitter Lesson.'' Future embodied foundation models might be trained through RL-like methods (such as tasks requiring environmental exploration) or other approaches to actively learn and rapidly adapt in new environments at inference, then perform new tasks effectively.

\Needspace{5\baselineskip}
\hypertarget{manifestations-of-embodied-ai-self-improvement-multiple-dimensions-and-environments}{%
\subsubsection{(2) Manifestations of Embodied AI Self-Improvement: Multiple Dimensions and Environments}\label{manifestations-of-embodied-ai-self-improvement-multiple-dimensions-and-environments}}

\Needspace{4\baselineskip}
\begin{enumerate}
\def\labelenumi{\arabic{enumi}.}
\tightlist
\item
  Self-improvement across multiple dimensions
\end{enumerate}

LLMs already show signs of participating in self-improvement, including model training and agent frameworks, but their interactive environments remain inside computers. Embodied AI's scope is broader, including model training, physical-structure design, manufacturing, testing, and iteration. It can acquire data autonomously and improve sensors and action modules to broaden channels for data collection and environmental interaction, expanding its perceptual and action spaces. From this perspective, inventing experimental techniques or measurement tools, discussed later, can also be viewed as part of embodied AI self-improvement.

\Needspace{4\baselineskip}
\begin{enumerate}
\def\labelenumi{\arabic{enumi}.}
\setcounter{enumi}{1}
\tightlist
\item
  Extending intelligence into more new environments
\end{enumerate}

Once embodied AI achieves autonomous learning and self-improvement, intelligence may extend into more environments. Agents can continually explore new environments, acquire data, and update understanding. Combining that understanding with self-improvement can produce more capable embodied systems for those environments, such as nanorobots, space-exploration robots, deep-sea robots, disaster-rescue robots, laboratory robots, and surgical robots. They can perceive, reason, and act in complex environments in real time instead of merely following preprogrammed paths. By collecting interaction data, training environment-specific models, manufacturing corresponding hardware, and improving data-acquisition methods, these robots can themselves iteratively develop stronger embodied intelligence for their environments. In this cycle, the ``perception--reasoning--action'' loop extends across scales, environments, and tasks.

\FloatBarrier
\Needspace{5\baselineskip}
\hypertarget{references-166}{%
\subsection{References}\label{references-166}}

\vspace{0.25em}
\begin{itemize}
\tightlist
\item
  Sutton, R. S. (2019). \href{http://www.incompleteideas.net/IncIdeas/BitterLesson.html}{The Bitter Lesson}. Incomplete Ideas.
\item
  Ha, D., \& Schmidhuber, J. (2018). \href{https://arxiv.org/abs/1803.10122}{World Models}. arXiv:1803.10122.
\item
  Physical Intelligence. (2026). \href{https://arxiv.org/abs/2604.15483}{π0.7: A Steerable Generalist Robotic Foundation Model with Emergent Capabilities}. arXiv:2604.15483.
\item
  Fan, J. (2026). \href{https://www.youtube.com/watch?v=3Y8aq_ofEVs}{Robotics' End Game: Nvidia's Jim Fan}. Sequoia Capital AI Ascent 2026 {[}Video{]}.
\end{itemize}

\clearpage
